\section{Threats to Validity and Limitations}
\label{sec:limits}

\paragraph{Population and ground-truth noise.}
Submissions to a top venue are self-selected above average, so the reject class
here is ``weak \emph{relative to the venue bar},'' not weak in absolute terms; we
do not generalize beyond this population. The accept/reject decision is itself
noisy and partly reflects reviewer disagreement, area-chair effects, and venue
fit \citep{cortes2021inconsistency,beygelzimer2023arbitrary}. We mitigate the
latter by also validating against the continuous mean reviewer rating (H4), but
no analysis here can exceed the reliability of its ground truth.

\paragraph{Leakage and contamination.}
Beyond verbatim memorization---excluded by construction through the post-cutoff
cohort---the model holds priors on institutions, authors, and topics that could
correlate with outcomes. We address this with the arXiv-before-cutoff sensitivity
split, the contaminated-contrast, and prospective checks detailed in
\S\ref{sec:discussion}, and the question of whether printed identity moves the
score by the separate prestige-perturbation experiment. Residual indirect
familiarity cannot be fully excluded retrospectively; only the prospective cohort
closes it entirely.

\paragraph{Ground truth may itself be partly machine-generated.}
A growing fraction of submitted peer reviews are now LLM-assisted at the venues we
study \citep{liang2024monitoring}. To the extent the reviewer ratings we validate
against are themselves shaped by language models, our agreement metric partly
measures agreement with other AI systems rather than with unaided human judgment.
We cannot separate the two components with the available data; the effect, if
present, inflates apparent agreement and is a ceiling on the interpretation of H4.

\paragraph{The cheap-model proxy.}
The scale of the headline rests on the bridge (H5), which we measure and gate on
rather than assume; had it failed, we would have led with the smaller frontier
cohort ($n=\NfullPrimary$) and wider intervals. The frontier cohort is modest by
budget and is used for confirmation, not as the sole evidence. The bridge is also
strongest globally and weaker exactly where the deployable claim lives: within the
full-mini bottom quintile the two scores agree at only $\rho=\bridgeLowRhoFull$ and
\bridgeOverlap\% of that band recurs in the frontier bottom quintile
(\S\ref{sec:res-validity}). The cheap-model score is therefore a valid large-$N$
proxy for the relationship, but a low cheap-model flag is best read as a screen to
be confirmed on the frontier score, not an exact substitute for it on an individual
submission. This is exactly why we report the deployable bottom-flag (H1) on the
frontier model directly (\S\ref{sec:res-validity},
\bottomRejectRateFrontier\%): the production claim does not route through the weak
low-end bridge, which constrains only the use of the cheap model as a standalone
bottom screen, not the headline.

\paragraph{An uninformative citation dimension.}
The citation subscore is pinned at the maximum on \citationPinnedFull\% of frontier
papers and discriminates at chance (AUROC \citationAUROC; \S\ref{sec:discussion}),
because the citation audit returned no recommendations for the original frontier
run and the pipeline scores an empty audit as a perfect bibliography---a silent
fallback, confirmed against re-grades that produce real scores. It does not affect
the headline---citation is the most lightly weighted dimension, and dropping it from
the weighting leaves discrimination unchanged (AUROC \aurocDropCitationFull{} vs.\
\aurocFull)---and the search-backed fix (and the lesson that grounded retrieval
enhances the model) is discussed in \S\ref{sec:discussion}.

\paragraph{Single venue and field.}
The primary cohort is one venue in one field, and we do not claim cross-field
generalization here. A planned follow-up grades a full review set at a second,
different-field venue; because venues that publish \emph{rejected} submissions are
rare, that study validates against the continuous reviewer rating---which
generalizes across fields---rather than the reject class, and is reported
separately. The within-family contaminated contrast (Appendix~\ref{app:robust})
addresses year-to-year stability only.

\paragraph{Defined semantics, fixed weights.}
A score's meaning is fixed by the grading rubric---each dimension is defined by
explicit criteria---so it is an anchored construct, and the overall is a fixed
weighted mean (Eq.~\ref{eq:overall}) set on substantive grounds and deliberately
\emph{not} fit to the outcome. Holding the weights label-independent is what
licenses the validity test. We validate that submissions \emph{ordered} by the
rubric track the human signal (discrimination and rank agreement); we do
\emph{not} calibrate the absolute level to an acceptance probability---a distinct,
venue-specific exercise we leave to future work. We report the dimension scores
and the weighting so a reader can reweight, but we do not optimize against the
decision.

\paragraph{Reproducibility with a proprietary scorer.}
The scored dataset, the decision and rating labels, the analysis code, and this
paper are released, so every figure, statistic, and the entire score-to-outcome
analysis reproduces independently. The grading rubric, prompts, and model
configuration are proprietary: the manuscript-to-score function is audited through
its outputs---and can be exercised on new manuscripts via the system---rather than
re-implemented. This is a defensible posture for validating a deployed instrument,
but weaker than full independent reproducibility of the scorer itself: the validation
is open even though the instrument is closed. Because the design was
frozen and pre-registered before unblinding, this openness is what lets a single
author's validation be independently checked rather than taken on trust. Additional
minor threats (post-rebuttal rating drift, citation relevance vs.\ groundedness,
run-to-run variance) are detailed in Appendix~\ref{app:limits}.

\paragraph{Sample-vs-eligible representativeness and per-review sensitivities.}
The population-boundary ledger (Appendix~\ref{app:robust},
Table~\ref{tab:popboundary}) reports counts and exclusion reasons for every
eligible submission, but a covariate-level comparison of the graded cohort to the
full eligible population is not possible because manuscript covariates (page and
word counts, reference and figure counts, primary area) were not extracted for
the ungraded submissions. Separately, the released two-CSV data carries only the
mean reviewer rating, so per-review sensitivities---a human-review reliability
ceiling and alternate rating aggregations (median or trimmed mean)---are deferred
to future work and available on request rather than reported here; we do not
synthesize per-review ratings to fill the gap.

\paragraph{What we do not claim.}
We do not claim the score predicts acceptance probability, ranks strong papers,
or should be used for positive selection or any consequential decision without a
human in the loop. The validated use is narrow and stated as such---the action we
validate is prioritization for human attention, not rejection. The
deployment-boundary, independence, and proprietary-scorer audit commitments this
posture rests on are collected in Appendix~\ref{app:posture}.
